\conclusion{Дүгнэлт}
Энэхүү судалгааны ажлын хүрээнд хоёр болон түүнээс дээш хүний хөдөлгөөнийг агуулсан видео бичлэгээс хүний биеийн гол цэгүүдийг тодорхойлж, тэдгээрийн хөдөлгөөний синхрончлолыг 0--100 хуваарьт тоон утгаар үнэлэх систем боловсруулж, хяналттай болон бодит бичлэгүүд дээр туршсан. Судалгааны зорилтуудын хүрээнд дараах үр дүнд хүрэв.

\begin{enumerate}
\def\labelenumi{\arabic{enumi}.}
\item
  Хүний байрлал тодорхойлох технологи, гол цэгийн форматууд болон хөдөлгөөний ижил төстэй байдлыг хэмжих аргуудыг судалж, үений өнцөгт суурилсан, хүний байрлал ба хэмжээнээс хамааралгүй дүрслэлийг сонгосон (Бүлэг 2).
\item
  Таван HPE загварыг ижил бичлэг дээр илрүүлэлт, гол цэгийн тохирол, чичиргээ, мөчдийн тогтвортой байдал, хурдаар харьцуулж, олон хүнийг тогтвортой илрүүлж, гол цэгийн чанараар хамгийн сайн үр дүн үзүүлсэн RTMPose-m загварыг сонгосон (3.2).
\item
  Хяналттай өгөгдлийг FFmpeg ашиглан нэг хүний бичлэгээс автоматаар үүсгэх арга боловсруулж, хоцрогдол, толин тусгал, өөр хөдөлгөөн, камерын зай, нягтаршил, кадрын давтамж, хүний тооны нийт 32 бичлэг үүсгэсэн (4.1).
\item
  Хүн хуваарилалт, дутуу цэгийг нөхөх, Савицкий-Голейн шүүлтүүрээр тэгшитгэх урьдчилсан боловсруулалтыг хэрэгжүүлсэн (3.3, 3.4).
\item
  Үений өнцгийн зөрүү, Пирсоны корреляц, хөндлөн корреляцаар олсон хоцрогдлыг нэгтгэсэн, кадр бүрийн цонх, биеийн хэсгийн жин, толин тусгал горим бүхий синхрончлолын алгоритм болон синхрончлол алдагдсан мөч, хоцорсон хүн, мөчийг илрүүлэх аргыг боловсруулсан (3.5, 3.6). Системийг вэб хөтөч дээр болон Python сервер дээр ажиллах хоёр горимтойгоор хэрэгжүүлсэн (3.7).
\item
  Видеоны түвшний хяналттай туршилтад алгоритм 0--1 секундийн бүх хоцрогдлыг 0 мс алдаатай олсон, оноо нь хоцрогдол нэмэгдэх тусам 99.8-аас 54.3 хүртэл жигд буурсан, нөхцлүүдийг синхрон \(>\) хоцрогдолтой \(>\) өөр хөдөлгөөн гэсэн дарааллаар ялгасан. Толин тусгал хөдөлгөөнийг 99.3 оноогоор зөв таньж, 3 ба 4 хүнтэй бичлэгт хоцорсон хүнийг зөв илрүүлсэн. Гол цэгийн түвшинд хоцроосон 20 тохиолдлын 18-д хоцорсон хүн болон мөчийг зөв тодорхойлсон. Сервер вэб хөтчөөс 2.2 дахин хурдан ажиллаж, хоёр горимын гол цэг медианаар 0.61 пиксел зөрүүтэй байв (4.3--4.5).
\item
  Хүний хэмжээг 50\%, нягтаршлыг 360p, кадрын давтамжийг 10 кадр/с хүртэл бууруулахад оноо болон хоцрогдлын үр дүн бараг өөрчлөгдөөгүй бол хүний өндөр ойролцоогоор 140 пикселээс, нягтаршил 240p-ээс бага үед хүний тоог автоматаар тодорхойлох алхам алдаа гаргасныг тогтоож, аргын боломж ба хязгаарлалтыг тодорхойлсон (4.4, 4.6).
\end{enumerate}

Мөн судалгааны явцад вэб хөтөч видеог кадрын хил дээр шилжүүлэхэд кадр давхардах, алгасагдах алдааг илрүүлж засварласан бөгөөд хоцрогдлын оноо \(\rho_{min}\) босгын орчимд тогтворгүй болохыг тогтоож, тогтворжуулах налуу хувилбарыг санал болгосон.

Иймд боловсруулсан систем нь энгийн нэг камераар авсан бичлэгээс олон хүний хөдөлгөөний синхрончлолыг тоон, давтагдах боломжтой хэлбэрээр үнэлж, хэн, хэзээ, биеийн аль хэсгээрээ хэр хэмжээгээр зөрж байгааг тодорхойлох боломжтой гэж дүгнэв.
